\documentclass[11pt]{article}
\usepackage{explorer}
\newcommand{\ExplorerLanguage}{FR}
\newcommand{\ExplorerNumber}{03}
\newcommand{\ExplorerKind}{Solutions détaillées}
\begin{document}
\ExplorerTitle{Lycée · arithmétique}{Congruences et derniers chiffres}
\ExplorerSection{Exercice 1}
Trouver le reste de −37 modulo 8.\par
−\(37 = 8 \times  (-5) + 3\), donc le reste est 3.\par
\ExplorerSection{Exercice 2}
Déterminer le dernier chiffre de \(7^{2026}\).\par
Le cycle est 7,9,3,1 et \(2026 \equiv  2\) (mod 4). Le dernier chiffre est 9.\par
\ExplorerSection{Exercice 3}
Prouver qu’un carré entier n’est jamais congru à 2 modulo 3.\par
Un entier est congru à 0,1 ou 2. Les carrés correspondants sont 0,1 et \(4 \equiv  1\) modulo 3.\par
\ExplorerSection{Exercice 4}
Résoudre \(3x \equiv  1\) (mod 7).\par
\(3 \times  5 = 15 \equiv  1\). Multiplier par 5 donne \(x \equiv  5\) (mod 7). Réciproquement, toutes ces valeurs conviennent.\par
\ExplorerSection{Problème : raisonnement complet}
Un entier N est congru à 2 modulo 3 et à 3 modulo 5. Écrire \(N = 3k + 2\), déterminer k modulo 5, puis donner toutes les solutions modulo 15. Parmi les solutions entre 0 et 100 inclus, lesquelles sont paires ?\par
\(3k + 2 \equiv  3\) (mod 5), donc \(3k \equiv  1\) et \(k \equiv  2\) (mod 5), car \(3 \times  2 \equiv  1\). Ainsi \(N = 15t + 8\). Dans [0,100], on obtient 8,23,38,53,68,83,98. Les valeurs paires sont 8,38,68,98. La substitution vérifie les deux congruences pour toute valeur entière de t.\par
\end{document}
